\documentclass[a4paper,11pt]{article}
\usepackage[T2A]{fontenc}
\usepackage[utf8]{inputenc}
\usepackage[russian]{babel}
\usepackage{amsmath,amssymb,mathtools}
\renewcommand{\familydefault}{\sfdefault}
\usepackage{geometry}
\geometry{left=23mm,right=23mm,top=18mm,bottom=20mm}
\usepackage{microtype}
\usepackage{tabularx,booktabs,array}
\usepackage{tikz}
\usepackage{xcolor}
\usepackage{hyperref}
\usepackage{fancyhdr}
\usepackage{enumitem}
\definecolor{accent}{HTML}{007D82}
\definecolor{ink}{HTML}{292D32}
\definecolor{wrong}{HTML}{B43D3D}
\hypersetup{colorlinks=true,linkcolor=accent}
\setlength{\parindent}{0pt}
\setlength{\parskip}{4pt}
\setlength{\abovedisplayskip}{7pt}
\setlength{\belowdisplayskip}{7pt}
\setlength{\abovedisplayshortskip}{4pt}
\setlength{\belowdisplayshortskip}{4pt}
\setlength{\headheight}{14pt}
\setlength{\emergencystretch}{2em}
\pagestyle{fancy}
\fancyhf{}
\fancyhead[L]{\small\color{ink}Ревью домашней работы}
\fancyhead[R]{\small\color{ink}Марк Гукин}
\fancyfoot[C]{\parbox{\linewidth}{\centering\scriptsize\color{ink}Лёвин Артём Александрович эксклюзивно для Марка Гукина\\[-1pt]\thepage}}
\renewcommand{\headrulewidth}{0.35pt}
\renewcommand{\footrulewidth}{0pt}
\newcommand{\keyidea}{\textcolor{blue!60!black}{\textbf{Ключевая идея.}}}
\newcommand{\mistake}{\textcolor{wrong}{\textbf{Ошибка.}}}
\newcommand{\result}{\textcolor{teal!60!black}{\textbf{Итог.}}}
\begin{document}
\thispagestyle{empty}
\begin{titlepage}
\centering
\vspace*{46mm}
{\color{accent}\large МАТЕМАТИКА\par}
\vspace{22mm}
{\fontsize{27}{32}\selectfont Ревью домашней работы\par}
\vfill
{\large Лёвин Артём Александрович\par}
\vspace{13mm}
{\large 4 ноября 2026 г.\par}
\vspace{21mm}
\begin{tikzpicture}
\draw[accent,line width=.8pt] (0,0) .. controls (1.15,.65) and (2.4,-.45) .. (3.5,.04)
  .. controls (4.65,.57) and (5.45,-.35) .. (6.65,.08)
  .. controls (7.85,.47) and (8.55,-.25) .. (10,0);
\end{tikzpicture}
\vspace{8mm}
\end{titlepage}

\newpage
\section*{Сводка по заданиям, требующим разбора}
\vspace{3mm}
\renewcommand{\arraystretch}{1.4}
\small
\begin{tabularx}{\linewidth}{@{}p{12mm}p{25mm}X p{21mm}p{24mm}@{}}
\toprule
№ & Тема & Суть ошибки & Ваш ответ & Верный ответ\\
\midrule
\hyperref[task:6]{\textbf{6}} & Формулы двойного угла & Ошибка при приведении косинуса угла $134^\circ$: потеряна функция косинуса & Не указан & $-18$\\
\bottomrule
\end{tabularx}
\normalsize
\vfill
\noindent\textcolor{accent}{\rule{34mm}{0.5pt}}

\newpage
\thispagestyle{empty}
\vspace*{0.34\textheight}
\begin{center}
{\fontsize{27}{33}\selectfont\color{accent}Разбор ошибок}
\end{center}
\vfill

\newpage
\phantomsection\label{task:6}
\section*{Задание № 6. Формулы двойного угла}

\textbf{Условие.} Вычислите:
\[
\frac{18\bigl(\sin^2 23^\circ-\cos^2 23^\circ\bigr)}{\cos46^\circ}.
\]

\textbf{Ваше решение (читаемая часть).}
\begin{align*}
&\frac{18\bigl(\sin^2 23^\circ-\cos^2 23^\circ\bigr)}{\cos46^\circ}\\
&=\frac{18\bigl(\cos^2 67^\circ-\sin^2 67^\circ\bigr)}{\cos46^\circ}
=\frac{18\cos134^\circ}{\cos46^\circ}
\end{align*}
Далее записано $\cos134^\circ=-\sin46^\circ$. Вычисление до числового ответа не доведено.

\textbf{Ваш ответ.} Не указан.

\mistake\ Вместо формулы приведения для косинуса использована формула с синусом.

\textbf{Где именно возникает ошибка.} Последний верный шаг:
\[
\frac{18\cos134^\circ}{\cos46^\circ}.
\]
Первый неверный переход:
\[
\textcolor{wrong}{\cos134^\circ=-\sin46^\circ}.
\]

\textbf{Почему это неверно.} Угол $134^\circ$ находится во II четверти, а
$134^\circ=180^\circ-46^\circ$. Для таких углов действует формула
\[
\cos(180^\circ-\alpha)=-\cos\alpha.
\]
Косинус меняет знак, но название функции сохраняется. Замена косинуса синусом здесь не допускается.

\keyidea\ В выражении встречаются одновременно углы $23^\circ$ и
$46^\circ$, причём $46^\circ=2\cdot23^\circ$. Поэтому удобно применить формулу двойного угла:
\[
\cos2\alpha=\cos^2\alpha-\sin^2\alpha.
\]

\textbf{Исправление от последнего правильного шага.}
\begin{align*}
\frac{18\cos134^\circ}{\cos46^\circ}
&=\frac{18\cos(180^\circ-46^\circ)}{\cos46^\circ}\\
&=\frac{-18\cos46^\circ}{\cos46^\circ}=-18.
\end{align*}

\textbf{Полное правильное решение (короткий маршрут).}
Из формулы двойного угла следует
\[
\sin^2 23^\circ-\cos^2 23^\circ
=-\bigl(\cos^2 23^\circ-\sin^2 23^\circ\bigr)
=-\cos46^\circ.
\]
Подставляем в исходное выражение:
\[
\frac{18(-\cos46^\circ)}{\cos46^\circ}=-18.
\]

\textbf{Проверка.} $\cos46^\circ\ne0$, поэтому сокращение допустимо.
Полученный числовой результат согласуется с отрицательным знаком
числителя и положительным знаменателем.

\result\ Ответ: $\boxed{-18}$.

\textbf{Задача на закрепление.} Вычислите:
\[
\frac{24\bigl(\sin^2 31^\circ-\cos^2 31^\circ\bigr)}{\cos62^\circ}.
\]
\end{document}
